\subsection{晶胞 string 关联与组合诊断量}
\label{sec:heatmaps}

价键链的 string 序~\cite{dennijs1989preroughening} 启发了与本文两格点晶胞匹配的非局域关联。
记 $M=L/2$、$T_r^z=(Z_{2r-1}+Z_{2r})/2$，定义
\begin{equation}
C_{\rm str}=
\left\langle T_1^z
\exp\left(i\pi\sum_{r=2}^{M-1}T_r^z\right)T_M^z\right\rangle.
\label{eq:cell-string}
\end{equation}
晶胞自旋 $T_r^z$ 的本征值为整数。
利用 $e^{i\pi T_r^z}=-Z_{2r-1}Z_{2r}$，得到计算和实验中实际使用的观测量：
\begin{align}
O_{\rm str}&=4C_{\rm str}\nonumber\\
&=\left\langle (Z_1+Z_2)
\prod_{r=2}^{M-1}(-Z_{2r-1}Z_{2r})
(Z_{L-1}+Z_L)\right\rangle.
\label{eq:measured-string}
\end{align}
$L=8$ 时两个内部负号抵消。系数四来自两个端点自旋之和。

string 序的对称性解释要求串内插入对称操作，且端点具有适当的变换性质~\cite{perezgarcia2008string,pollmann2012detection}。
本文 $u_xT_r^zu_x^\dagger=-T_r^z$，因此端点在第二个晶胞对称操作下为奇。当端点位于对称、可注入 MPS 的体内时，长 $u_z$ 串的选择规则要求该关联在 $\eta=+1$ 时消失，$\eta=-1$ 时则允许非零。有限开链的首尾关联同时包含体与边界态的贡献。
包含交换操作的对称性不变量与 string 序诊断互为补充~\cite{haegeman2012order}。
相关非局域关联已在 Mott 绝缘体~\cite{endres2011observation} 和 SPT 自旋链~\cite{tan2023realizing} 中测量。
本文式~\eqref{eq:measured-string} 是有限长度诊断量。

数值热力图采用固定尺度
\begin{equation}
\hat S(\pi)=\frac{S(\pi)}{L},\qquad
\hat O_{\rm str}=-O_{\rm str}.
\label{eq:fixed-scales}
\end{equation}
结构因子满足 $0\leq\hat S(\pi)\leq1$；string 的符号约定使拓扑参考态满足 $\hat O_{\rm str}=1$。

在本模型的基态族中，两种关联量提供互补信息：交错结构因子识别 AFM 区域，晶胞 string 关联则反映保持对称性的平庸与拓扑二聚化区域之间的差别。
将二者组合成单一诊断量
\begin{equation}
Q=\frac43-2\hat O_{\rm str}-\frac43\hat S(\pi).
\label{eq:Q}
\end{equation}
其系数可直接从三个 $L=8$ 参考态理解：奇键单态直积态的 $(\hat S,\hat O)=(1/4,0)$，体单态加首尾单态的参考态为 $(1/4,1)$，AFM 猫态为 $(1,0)$。
式~\eqref{eq:Q} 是将三者依次赋值为 $Q=+1,-1,0$ 的唯一仿射组合。
这些参考值由态本身给出，每个单态均贡献交错反关联，并共同确定 $Q$ 的参考态尺度。

图~\ref{fig:D04}(a)--(c) 给出与反射标记在同一网格上计算的 $\hat S(\pi)$、$\hat O_{\rm str}$ 和 $Q$。
共同网格使各量的参数依赖可以直接比较。
组合诊断量与部分反射标记呈现相近的三相区域结构，但有限尺寸下的转变轮廓有所不同。
$L=8$ 时，式~\eqref{eq:Q} 等价于 $Q=4/3+2O_{\rm str}-S(\pi)/6$。

实验扫描选择 $Q$，是因为两种关联量在计算基均为对角量：同一组比特串即可估计两者的联合取值和协方差，每个态制备电路只需一种测量基。
对一般制备态，仅凭这些计算基概率无法确定 $\tilde Z_{\mathcal R}$，后者还依赖镜像自旋组态之间的相干信息和子系统纯度。
文献~\cite{elben2020manybody} 的随机测量方案通过多个测量设置下镜像关联的局域随机旋转，估计反射重叠及归一化所需的纯度。
采用 $Q$ 可省去这些额外的测量基设置，为研究本模型的三个区域及其转变提供可直接实验获取的替代诊断。
它是经参考态标定的关联诊断量，而非量子化拓扑不变量；$\tilde Z_{\mathcal R}$ 提供基于对称性的数值参照。

相关方法还包括用于预测观测量的经典阴影~\cite{huang2020predicting}，以及用于相识别的量子卷积电路~\cite{cong2019quantum}。
第~\ref{sec:comparison} 节说明实验数据另行拟合的尺度。
